\section*{September 8, 2026: Candidate geography and boundary distances}
\textit{Agent-drafted reconstruction and independent spatial checks.}

The candidate geography now builds through construction-year boundary assignment.
It accounts for 14,017 candidate project-years: 13,903 have complete parcel shapes
and distances, and 114 remain geographically unresolved. Another 23 candidates in
the review queue have no selected construction year. The located candidates include
4,214 within 500 feet of a ward boundary. These are preliminary groups from the
Assessor and commercial records, not the paper's final estimation sample; duplicate
identity, measurements, and eligibility still require downstream reconciliation.

The production scripts combine a project's component parcels only when every
component has an accepted match. They calculate the centroid of the combined
shape and assign the ward map using June 15 of the recorded construction year,
following the convention already stated in the paper. Incomplete projects retain
missing wards and distances. Additional assertions prevent duplicate parcel keys,
accepted components absent from the request table, nonpositive areas, and ambiguous
ward membership from passing silently.

An independent calculation compares each centroid with every shared boundary of
its assigned ward, rather than using the production nearest-feature helper. All
13,903 saved distances and chosen ward pairs agree within 0.000001 feet. This
verifies the distance calculation for the supplied project shapes; it does not
establish that every source parcel belongs to the claimed project.

That distinction matters for 41 centroids outside their assembled parcel shapes,
including nine within 500 feet of a ward boundary. Most gaps are small and can
arise from irregular or disconnected parcels. The largest instead reflects an
Assessor tieback linking parcels 20174140320000 and 29174140310000 about 12.6 miles
apart. The candidate was already marked for review because no complete measurement
snapshot could be reconciled. Its calculated center should not be interpreted as
a verified building location. These flags remain review items, without automatic
exclusions or replacement coordinates.

One consolidated list contains 194 review items across 188 project identifiers.
It combines unresolved component locations, the 23 missing years, the centroid
flags, and all 30 earlier predecessor checkpoints, accounting for overlap. The
one-year-ahead exact-parcel coordinates for 763 W 15th and 3609 W 50th remain
recorded evidence; they have not been adopted. The list is generated from the
current outputs, without case-specific conditions in the code.

The annual timing assumption is also consequential. Among 754 located candidates
dated 2015, using the old map instead changes 179 ward assignments and 189
indicators for being within 500 feet. This sensitivity calculation does not change
the paper's June 15 convention or establish that an individual date is wrong.

\textit{Reproduction and limits.} Work follows \path{a5b89993} on
\path{sales_sample_exploration}. Build the six preferred geography and boundary
report targets in \path{tasks/new_construction_cleaning/code/}, then the
\path{candidate_geography_checks.csv} and \path{candidate_geography_review_queue.csv}
reports in the release audit. The pinned parcel inputs and existing ward maps feed
these targets through Make. All six outputs reproduce their pre-assertion contents; each missing output also
regenerates once under parallel Make, and an unchanged build runs no producers.
The paper input and estimates remain unchanged.
Logbook compilation retains the documented holds on unrelated unfinished stages.
This completes the candidate-geography milestone, not the full replication or
the final density dataset.
